\documentclass[11pt]{article}
% =============================================================================
%  Shared article preamble — tutorials, homework, solutions.
%  Same fonts, palette and notation as the Beamer decks (see preamble.tex).
%      \documentclass[11pt]{article}
%      % =============================================================================
%  Shared article preamble — tutorials, homework, solutions.
%  Same fonts, palette and notation as the Beamer decks (see preamble.tex).
%      \documentclass[11pt]{article}
%      % =============================================================================
%  Shared article preamble — tutorials, homework, solutions.
%  Same fonts, palette and notation as the Beamer decks (see preamble.tex).
%      \documentclass[11pt]{article}
%      \input{../../../style/preamble_article}
% =============================================================================
\usepackage[margin=1in]{geometry}
\usepackage{helvet}
\usepackage[default]{lato}
\usepackage{mathpazo}
\usepackage[utf8]{inputenc}
\usepackage{amsmath,amssymb,amsthm}
\usepackage{mathtools}
\usepackage{graphicx}
\usepackage{booktabs}
\usepackage{multirow}
\usepackage{enumitem}
\usepackage{xcolor}
\usepackage{tcolorbox}
\tcbuselibrary{skins,breakable}
\usepackage{listings}
\usepackage[ruled,vlined]{algorithm2e}
\usepackage{hyperref}
\usepackage{fancyhdr}
\usepackage{titlesec}

\definecolor{blue}{RGB}{0,114,178}
\definecolor{red}{RGB}{213,94,0}
\definecolor{yellow}{RGB}{240,228,66}
\definecolor{green}{RGB}{0,158,115}
\hypersetup{colorlinks=true, linkcolor=blue, urlcolor=blue, citecolor=blue}

\titleformat{\section}{\Large\bfseries\color{blue}}{\thesection}{0.6em}{}
\titleformat{\subsection}{\large\bfseries}{\thesubsection}{0.6em}{}
\pagestyle{fancy}\fancyhf{}
\rhead{\footnotesize ENEI -- INEI \,\textperiodcentered\, PEU-CD 2026}
\lfoot{\footnotesize Machine Learning I \& II \,\textperiodcentered\, Alexander Quispe}
\rfoot{\footnotesize\thepage}
\renewcommand{\headrulewidth}{0.3pt}

\newtcolorbox{derivation}[1]{enhanced, breakable, colback=blue!3, colframe=blue, fonttitle=\bfseries,
  title={#1}, boxrule=0.6pt, left=6pt, right=6pt, top=4pt, bottom=4pt, arc=2pt}
\newtcolorbox{claim}[1]{enhanced, breakable, colback=white, colframe=black, fonttitle=\bfseries,
  title={#1}, boxrule=0.6pt, left=6pt, right=6pt, top=4pt, bottom=4pt, arc=2pt}
\newtcolorbox{keypoint}{enhanced, breakable, colback=green!6, colframe=green, boxrule=0.6pt,
  left=6pt, right=6pt, top=4pt, bottom=4pt, arc=2pt}
\newtcolorbox{caution}{enhanced, breakable, colback=red!5, colframe=red, boxrule=0.6pt,
  left=6pt, right=6pt, top=4pt, bottom=4pt, arc=2pt}
\newtcolorbox{task}[1]{enhanced, breakable, colback=yellow!12, colframe=yellow!60!black, fonttitle=\bfseries,
  title={#1}, boxrule=0.6pt, left=6pt, right=6pt, top=4pt, bottom=4pt, arc=2pt}

\lstset{
  basicstyle=\ttfamily\small, keywordstyle=\color{blue}\bfseries,
  commentstyle=\color{green}, stringstyle=\color{red},
  showstringspaces=false, breaklines=true, frame=single, rulecolor=\color{black!30},
  numbers=none, xleftmargin=6pt, framexleftmargin=6pt, language=Python
}
\newtheorem{theorem}{Theorem}
\newtheorem{lemma}{Lemma}
\newtheorem{proposition}{Proposition}
\theoremstyle{definition}
\newtheorem{definition}{Definition}
\newtheorem{exercise}{Exercise}

\input{../../../style/macros}

% =============================================================================
\usepackage[margin=1in]{geometry}
\usepackage{helvet}
\usepackage[default]{lato}
\usepackage{mathpazo}
\usepackage[utf8]{inputenc}
\usepackage{amsmath,amssymb,amsthm}
\usepackage{mathtools}
\usepackage{graphicx}
\usepackage{booktabs}
\usepackage{multirow}
\usepackage{enumitem}
\usepackage{xcolor}
\usepackage{tcolorbox}
\tcbuselibrary{skins,breakable}
\usepackage{listings}
\usepackage[ruled,vlined]{algorithm2e}
\usepackage{hyperref}
\usepackage{fancyhdr}
\usepackage{titlesec}

\definecolor{blue}{RGB}{0,114,178}
\definecolor{red}{RGB}{213,94,0}
\definecolor{yellow}{RGB}{240,228,66}
\definecolor{green}{RGB}{0,158,115}
\hypersetup{colorlinks=true, linkcolor=blue, urlcolor=blue, citecolor=blue}

\titleformat{\section}{\Large\bfseries\color{blue}}{\thesection}{0.6em}{}
\titleformat{\subsection}{\large\bfseries}{\thesubsection}{0.6em}{}
\pagestyle{fancy}\fancyhf{}
\rhead{\footnotesize ENEI -- INEI \,\textperiodcentered\, PEU-CD 2026}
\lfoot{\footnotesize Machine Learning I \& II \,\textperiodcentered\, Alexander Quispe}
\rfoot{\footnotesize\thepage}
\renewcommand{\headrulewidth}{0.3pt}

\newtcolorbox{derivation}[1]{enhanced, breakable, colback=blue!3, colframe=blue, fonttitle=\bfseries,
  title={#1}, boxrule=0.6pt, left=6pt, right=6pt, top=4pt, bottom=4pt, arc=2pt}
\newtcolorbox{claim}[1]{enhanced, breakable, colback=white, colframe=black, fonttitle=\bfseries,
  title={#1}, boxrule=0.6pt, left=6pt, right=6pt, top=4pt, bottom=4pt, arc=2pt}
\newtcolorbox{keypoint}{enhanced, breakable, colback=green!6, colframe=green, boxrule=0.6pt,
  left=6pt, right=6pt, top=4pt, bottom=4pt, arc=2pt}
\newtcolorbox{caution}{enhanced, breakable, colback=red!5, colframe=red, boxrule=0.6pt,
  left=6pt, right=6pt, top=4pt, bottom=4pt, arc=2pt}
\newtcolorbox{task}[1]{enhanced, breakable, colback=yellow!12, colframe=yellow!60!black, fonttitle=\bfseries,
  title={#1}, boxrule=0.6pt, left=6pt, right=6pt, top=4pt, bottom=4pt, arc=2pt}

\lstset{
  basicstyle=\ttfamily\small, keywordstyle=\color{blue}\bfseries,
  commentstyle=\color{green}, stringstyle=\color{red},
  showstringspaces=false, breaklines=true, frame=single, rulecolor=\color{black!30},
  numbers=none, xleftmargin=6pt, framexleftmargin=6pt, language=Python
}
\newtheorem{theorem}{Theorem}
\newtheorem{lemma}{Lemma}
\newtheorem{proposition}{Proposition}
\theoremstyle{definition}
\newtheorem{definition}{Definition}
\newtheorem{exercise}{Exercise}

% =============================================================================
%  Notation — one convention for all lectures
%
%  scalars      x, y, w           plain italic
%  vectors      \vx \vy \vw       bold lower-case
%  matrices     \vX \vW \vSigma   bold upper-case
%  parameters   \vbeta \vtheta    bold Greek
%  transpose    \vX\T             \top
% =============================================================================
\newcommand{\R}{\mathbb{R}}
\newcommand{\E}{\mathbb{E}}
\newcommand{\Prob}{\mathbb{P}}
\DeclareMathOperator{\Var}{Var}
\DeclareMathOperator{\Cov}{Cov}
\DeclareMathOperator{\tr}{tr}
\DeclareMathOperator{\diag}{diag}
\DeclareMathOperator{\rank}{rank}
\DeclareMathOperator{\sign}{sign}
\DeclareMathOperator{\col}{col}
\DeclareMathOperator{\RSS}{RSS}
\DeclareMathOperator*{\argmin}{arg\,min}
\DeclareMathOperator*{\argmax}{arg\,max}
\newcommand{\norm}[1]{\left\lVert #1 \right\rVert}
\newcommand{\abs}[1]{\left\lvert #1 \right\rvert}
\newcommand{\inner}[2]{\left\langle #1,\, #2 \right\rangle}
\newcommand{\ind}[1]{\mathbf{1}\!\left\{#1\right\}}
\newcommand{\T}{^{\top}}
\newcommand{\inv}{^{-1}}
\newcommand{\half}{\tfrac{1}{2}}
\newcommand{\eps}{\varepsilon}
\newcommand{\Dset}{\mathcal{D}}
\newcommand{\Hset}{\mathcal{H}}
\newcommand{\Lcal}{\mathcal{L}}
\newcommand{\Ncal}{\mathcal{N}}

% bold vectors
\newcommand{\vx}{\mathbf{x}}   \newcommand{\vy}{\mathbf{y}}   \newcommand{\vz}{\mathbf{z}}
\newcommand{\vw}{\mathbf{w}}   \newcommand{\vu}{\mathbf{u}}   \newcommand{\vv}{\mathbf{v}}
\newcommand{\vb}{\mathbf{b}}   \newcommand{\vh}{\mathbf{h}}   \newcommand{\vr}{\mathbf{r}}
\newcommand{\vc}{\mathbf{c}}   \newcommand{\vf}{\mathbf{f}}   \newcommand{\vg}{\mathbf{g}}
\newcommand{\vp}{\mathbf{p}}   \newcommand{\vq}{\mathbf{q}}   \newcommand{\vs}{\mathbf{s}}
\newcommand{\va}{\mathbf{a}}   \newcommand{\vd}{\mathbf{d}}   \newcommand{\vt}{\mathbf{t}}
\newcommand{\vk}{\mathbf{k}}   \newcommand{\vm}{\mathbf{m}}   \newcommand{\vn}{\mathbf{n}}
\newcommand{\vzero}{\mathbf{0}} \newcommand{\vone}{\mathbf{1}}
% bold matrices
\newcommand{\vX}{\mathbf{X}}   \newcommand{\vY}{\mathbf{Y}}   \newcommand{\vZ}{\mathbf{Z}}
\newcommand{\vW}{\mathbf{W}}   \newcommand{\vH}{\mathbf{H}}   \newcommand{\vI}{\mathbf{I}}
\newcommand{\vA}{\mathbf{A}}   \newcommand{\vB}{\mathbf{B}}   \newcommand{\vC}{\mathbf{C}}
\newcommand{\vD}{\mathbf{D}}   \newcommand{\vP}{\mathbf{P}}   \newcommand{\vQ}{\mathbf{Q}}
\newcommand{\vU}{\mathbf{U}}   \newcommand{\vV}{\mathbf{V}}   \newcommand{\vS}{\mathbf{S}}
\newcommand{\vM}{\mathbf{M}}   \newcommand{\vK}{\mathbf{K}}   \newcommand{\vG}{\mathbf{G}}
% bold Greek
\newcommand{\vbeta}{\boldsymbol{\beta}}     \newcommand{\valpha}{\boldsymbol{\alpha}}
\newcommand{\vtheta}{\boldsymbol{\theta}}   \newcommand{\vmu}{\boldsymbol{\mu}}
\newcommand{\vSigma}{\boldsymbol{\Sigma}}   \newcommand{\vLambda}{\boldsymbol{\Lambda}}
\newcommand{\veps}{\boldsymbol{\varepsilon}} \newcommand{\vphi}{\boldsymbol{\phi}}
\newcommand{\vgamma}{\boldsymbol{\gamma}}   \newcommand{\vlambda}{\boldsymbol{\lambda}}
\newcommand{\vdelta}{\boldsymbol{\delta}}   \newcommand{\vnu}{\boldsymbol{\nu}}
\newcommand{\vxi}{\boldsymbol{\xi}}         \newcommand{\vpi}{\boldsymbol{\pi}}
\newcommand{\vomega}{\boldsymbol{\omega}}



% =============================================================================
\usepackage[margin=1in]{geometry}
\usepackage{helvet}
\usepackage[default]{lato}
\usepackage{mathpazo}
\usepackage[utf8]{inputenc}
\usepackage{amsmath,amssymb,amsthm}
\usepackage{mathtools}
\usepackage{graphicx}
\usepackage{booktabs}
\usepackage{multirow}
\usepackage{enumitem}
\usepackage{xcolor}
\usepackage{tcolorbox}
\tcbuselibrary{skins,breakable}
\usepackage{listings}
\usepackage[ruled,vlined]{algorithm2e}
\usepackage{hyperref}
\usepackage{fancyhdr}
\usepackage{titlesec}

\definecolor{blue}{RGB}{0,114,178}
\definecolor{red}{RGB}{213,94,0}
\definecolor{yellow}{RGB}{240,228,66}
\definecolor{green}{RGB}{0,158,115}
\hypersetup{colorlinks=true, linkcolor=blue, urlcolor=blue, citecolor=blue}

\titleformat{\section}{\Large\bfseries\color{blue}}{\thesection}{0.6em}{}
\titleformat{\subsection}{\large\bfseries}{\thesubsection}{0.6em}{}
\pagestyle{fancy}\fancyhf{}
\rhead{\footnotesize ENEI -- INEI \,\textperiodcentered\, PEU-CD 2026}
\lfoot{\footnotesize Machine Learning I \& II \,\textperiodcentered\, Alexander Quispe}
\rfoot{\footnotesize\thepage}
\renewcommand{\headrulewidth}{0.3pt}

\newtcolorbox{derivation}[1]{enhanced, breakable, colback=blue!3, colframe=blue, fonttitle=\bfseries,
  title={#1}, boxrule=0.6pt, left=6pt, right=6pt, top=4pt, bottom=4pt, arc=2pt}
\newtcolorbox{claim}[1]{enhanced, breakable, colback=white, colframe=black, fonttitle=\bfseries,
  title={#1}, boxrule=0.6pt, left=6pt, right=6pt, top=4pt, bottom=4pt, arc=2pt}
\newtcolorbox{keypoint}{enhanced, breakable, colback=green!6, colframe=green, boxrule=0.6pt,
  left=6pt, right=6pt, top=4pt, bottom=4pt, arc=2pt}
\newtcolorbox{caution}{enhanced, breakable, colback=red!5, colframe=red, boxrule=0.6pt,
  left=6pt, right=6pt, top=4pt, bottom=4pt, arc=2pt}
\newtcolorbox{task}[1]{enhanced, breakable, colback=yellow!12, colframe=yellow!60!black, fonttitle=\bfseries,
  title={#1}, boxrule=0.6pt, left=6pt, right=6pt, top=4pt, bottom=4pt, arc=2pt}

\lstset{
  basicstyle=\ttfamily\small, keywordstyle=\color{blue}\bfseries,
  commentstyle=\color{green}, stringstyle=\color{red},
  showstringspaces=false, breaklines=true, frame=single, rulecolor=\color{black!30},
  numbers=none, xleftmargin=6pt, framexleftmargin=6pt, language=Python
}
\newtheorem{theorem}{Theorem}
\newtheorem{lemma}{Lemma}
\newtheorem{proposition}{Proposition}
\theoremstyle{definition}
\newtheorem{definition}{Definition}
\newtheorem{exercise}{Exercise}

% =============================================================================
%  Notation — one convention for all lectures
%
%  scalars      x, y, w           plain italic
%  vectors      \vx \vy \vw       bold lower-case
%  matrices     \vX \vW \vSigma   bold upper-case
%  parameters   \vbeta \vtheta    bold Greek
%  transpose    \vX\T             \top
% =============================================================================
\newcommand{\R}{\mathbb{R}}
\newcommand{\E}{\mathbb{E}}
\newcommand{\Prob}{\mathbb{P}}
\DeclareMathOperator{\Var}{Var}
\DeclareMathOperator{\Cov}{Cov}
\DeclareMathOperator{\tr}{tr}
\DeclareMathOperator{\diag}{diag}
\DeclareMathOperator{\rank}{rank}
\DeclareMathOperator{\sign}{sign}
\DeclareMathOperator{\col}{col}
\DeclareMathOperator{\RSS}{RSS}
\DeclareMathOperator*{\argmin}{arg\,min}
\DeclareMathOperator*{\argmax}{arg\,max}
\newcommand{\norm}[1]{\left\lVert #1 \right\rVert}
\newcommand{\abs}[1]{\left\lvert #1 \right\rvert}
\newcommand{\inner}[2]{\left\langle #1,\, #2 \right\rangle}
\newcommand{\ind}[1]{\mathbf{1}\!\left\{#1\right\}}
\newcommand{\T}{^{\top}}
\newcommand{\inv}{^{-1}}
\newcommand{\half}{\tfrac{1}{2}}
\newcommand{\eps}{\varepsilon}
\newcommand{\Dset}{\mathcal{D}}
\newcommand{\Hset}{\mathcal{H}}
\newcommand{\Lcal}{\mathcal{L}}
\newcommand{\Ncal}{\mathcal{N}}

% bold vectors
\newcommand{\vx}{\mathbf{x}}   \newcommand{\vy}{\mathbf{y}}   \newcommand{\vz}{\mathbf{z}}
\newcommand{\vw}{\mathbf{w}}   \newcommand{\vu}{\mathbf{u}}   \newcommand{\vv}{\mathbf{v}}
\newcommand{\vb}{\mathbf{b}}   \newcommand{\vh}{\mathbf{h}}   \newcommand{\vr}{\mathbf{r}}
\newcommand{\vc}{\mathbf{c}}   \newcommand{\vf}{\mathbf{f}}   \newcommand{\vg}{\mathbf{g}}
\newcommand{\vp}{\mathbf{p}}   \newcommand{\vq}{\mathbf{q}}   \newcommand{\vs}{\mathbf{s}}
\newcommand{\va}{\mathbf{a}}   \newcommand{\vd}{\mathbf{d}}   \newcommand{\vt}{\mathbf{t}}
\newcommand{\vk}{\mathbf{k}}   \newcommand{\vm}{\mathbf{m}}   \newcommand{\vn}{\mathbf{n}}
\newcommand{\vzero}{\mathbf{0}} \newcommand{\vone}{\mathbf{1}}
% bold matrices
\newcommand{\vX}{\mathbf{X}}   \newcommand{\vY}{\mathbf{Y}}   \newcommand{\vZ}{\mathbf{Z}}
\newcommand{\vW}{\mathbf{W}}   \newcommand{\vH}{\mathbf{H}}   \newcommand{\vI}{\mathbf{I}}
\newcommand{\vA}{\mathbf{A}}   \newcommand{\vB}{\mathbf{B}}   \newcommand{\vC}{\mathbf{C}}
\newcommand{\vD}{\mathbf{D}}   \newcommand{\vP}{\mathbf{P}}   \newcommand{\vQ}{\mathbf{Q}}
\newcommand{\vU}{\mathbf{U}}   \newcommand{\vV}{\mathbf{V}}   \newcommand{\vS}{\mathbf{S}}
\newcommand{\vM}{\mathbf{M}}   \newcommand{\vK}{\mathbf{K}}   \newcommand{\vG}{\mathbf{G}}
% bold Greek
\newcommand{\vbeta}{\boldsymbol{\beta}}     \newcommand{\valpha}{\boldsymbol{\alpha}}
\newcommand{\vtheta}{\boldsymbol{\theta}}   \newcommand{\vmu}{\boldsymbol{\mu}}
\newcommand{\vSigma}{\boldsymbol{\Sigma}}   \newcommand{\vLambda}{\boldsymbol{\Lambda}}
\newcommand{\veps}{\boldsymbol{\varepsilon}} \newcommand{\vphi}{\boldsymbol{\phi}}
\newcommand{\vgamma}{\boldsymbol{\gamma}}   \newcommand{\vlambda}{\boldsymbol{\lambda}}
\newcommand{\vdelta}{\boldsymbol{\delta}}   \newcommand{\vnu}{\boldsymbol{\nu}}
\newcommand{\vxi}{\boldsymbol{\xi}}         \newcommand{\vpi}{\boldsymbol{\pi}}
\newcommand{\vomega}{\boldsymbol{\omega}}




\title{\textcolor{blue}{Lab 1 --- Working Environment and a First End-to-End Pipeline}\\[4pt]
\large Machine Learning I \,\textperiodcentered\, Tutorial}
\author{Rodrigo Grijalba (teaching assistant) \,\textperiodcentered\, Alexander Quispe (instructor)}
\date{Wednesday, September 9, 2026 \,\textperiodcentered\, 19:00--22:00}

\begin{document}
\maketitle

\begin{keypoint}
\textbf{Goal of the session.} Leave with a working Python environment, a copy of the course repository,
and a complete supervised-learning pipeline you wrote yourself: data $\to$ split $\to$ fit $\to$ validate
$\to$ report. Every number in Lecture 1's worked example gets recomputed on the way.
\end{keypoint}

\section*{Materials}
\begin{itemize}[nosep]
  \item \texttt{lab\_01\_student.ipynb} --- the notebook you work in. Cells marked \texttt{\# TODO} are yours.
  \item \texttt{lab\_01\_solution.ipynb} --- released after the session.
  \item Sources: CS 155 Recitation 1 (Python) and Problem Set 1 (bias--variance, SGD).
\end{itemize}

\section{Environment (20 min)}

\subsection{Option A --- Google Colab}
Open \url{https://colab.research.google.com}, upload \texttt{lab\_01\_student.ipynb}. Everything needed is
preinstalled. Runtime $\to$ Change runtime type $\to$ CPU is enough for ML I.

\subsection{Option B --- local}
\begin{lstlisting}[language=bash]
python3 -m venv .venv && source .venv/bin/activate      # Windows: .venv\Scripts\activate
pip install numpy pandas scikit-learn matplotlib jupyter
jupyter lab
\end{lstlisting}

\subsection{The repository}
\begin{lstlisting}[language=bash]
git lfs install                       # once; the reference PDFs are stored with LFS
git clone https://github.com/alexanderquispe/enei-machine-learning-2026.git
cd enei-machine-learning-2026/ml1/labs/lab_01_pipeline
\end{lstlisting}
Pull before each session: \texttt{git pull}. You will never need to push.

\section{numpy for Machine Learning (30 min)}

Everything in this course is a matrix computation. Five things to be fluent in:
\begin{enumerate}[nosep]
  \item \textbf{Shapes.} \texttt{X.shape} is \texttt{(N, p)}; \texttt{y.shape} is \texttt{(N,)}. A shape mismatch
        is the most common bug you will write this term.
  \item \textbf{Matrix product.} \texttt{X.T @ X}, never \texttt{X.T * X} (that is elementwise).
  \item \textbf{Broadcasting.} \texttt{X - X.mean(axis=0)} subtracts the column means from every row.
  \item \textbf{Solving, not inverting.} \texttt{np.linalg.solve(A, b)} rather than \texttt{np.linalg.inv(A) @ b}.
        Same answer, better numerics, and it fails loudly on a singular \texttt{A}.
  \item \textbf{Randomness.} \texttt{rng = np.random.default\_rng(155)}; then \texttt{rng.normal}, \texttt{rng.permutation}.
        Seeded, so a result can be reproduced.
\end{enumerate}

\section{Lecture 1, Recomputed (40 min)}

\subsection{The worked example}
\begin{task}{Task 1 --- normal equations by hand}
With $x=(1,2,3,4,5)$ and $y=(2.2,2.8,4.5,3.7,5.5)$, build the design matrix $\vX$ (column of ones first),
form $\vX\T\vX$ and $\vX\T\vy$, and solve. You must obtain $\hat\vbeta=(1.49,\,0.75)\T$ --- Lecture 1, slide 15.
\end{task}

\subsection{The hat matrix, numerically}
\begin{task}{Task 2 --- four properties of $\vH=\vX(\vX\T\vX)\inv\vX\T$}
Compute $\vH$ ($5\times5$) and verify, to floating-point precision:
\begin{enumerate}[nosep]
  \item $\vH\T=\vH$ \hfill (symmetric)
  \item $\vH\vH=\vH$ \hfill (idempotent)
  \item $\tr\vH=2=p+1$ \hfill (degrees of freedom)
  \item $\vX\T(\vy-\vH\vy)=\vzero$ \hfill (residuals orthogonal to the columns)
\end{enumerate}
Use \texttt{np.allclose}. Then check that the eigenvalues of $\vH$ are $\{1,1,0,0,0\}$ and explain why
(Lecture 1, slide 20).
\end{task}

\subsection{Gradient descent}
\begin{task}{Task 3 --- implement, then break it}
Implement gradient descent for $L(\vbeta)=\frac1{2N}\norm{\vy-\vX\vbeta}^2$, starting at $\vbeta=\vzero$.
\begin{enumerate}[nosep]
  \item With $\eta=0.05$ and 2000 iterations, confirm convergence to $(1.49,0.75)$.
  \item Compute $\lambda_{\max}(\vX\T\vX/N)$ and the critical step $2/\lambda_{\max}$. Run with $\eta$ slightly
        below and slightly above it. Plot $L$ against iteration in both cases.
  \item Standardize $x$ (subtract mean, divide by standard deviation) and repeat. How does $\lambda_{\max}$ change?
        How many iterations does convergence take now?
\end{enumerate}
\end{task}
The last item is the practical lesson of the section: \emph{scale your features}. It is the cheapest
conditioning fix there is.

\section{A Real Pipeline with scikit-learn (50 min)}

The dataset is \texttt{sklearn.datasets.load\_diabetes}: $N=442$ patients, $p=10$ standardized
covariates, a quantitative disease-progression score as target. It ships with scikit-learn; no download.

\subsection{The protocol}
\begin{lstlisting}
from sklearn.model_selection import train_test_split
X_tr, X_te, y_tr, y_te = train_test_split(X, y, test_size=0.2, random_state=155)
\end{lstlisting}
The test set is now sealed. It is opened once, at the very end, to report one number.

\subsection{Baseline and first model}
\begin{task}{Task 4 --- two numbers to beat}
\begin{enumerate}[nosep]
  \item Baseline: predict $\bar y_{\text{train}}$ for everyone. Compute its MSE on the training set.
        This is the number any model must beat.
  \item Fit \texttt{LinearRegression} on the training set. Report training MSE and $R^2$.
        Verify that \texttt{model.coef\_} equals what \texttt{np.linalg.solve} gives on the same data.
\end{enumerate}
\end{task}

\subsection{$K$-fold cross-validation, by hand}
\begin{task}{Task 5 --- write \texttt{kfold\_mse}}
Write a function that takes $(\vX,\vy,K,\text{model})$, splits the rows into $K$ folds with a seeded
permutation, trains on $K-1$ folds and evaluates on the held-out one, and returns the $K$ validation MSEs.
Compare its mean with \texttt{sklearn.model\_selection.cross\_val\_score(..., scoring="neg\_mean\_squared\_error")}.
They should agree when both use the same fold assignment (\texttt{KFold(K, shuffle=True, random\_state=155)}).
\end{task}

\subsection{Model selection: the bias--variance picture, for real}
\begin{task}{Task 6 --- polynomial degree}
Using \texttt{PolynomialFeatures(degree=d)} for $d=1,2,3$, compute training MSE and 5-fold CV MSE for each.
Plot both against $d$. Training error must fall monotonically; CV error should rise by $d=2$ or $3$.
Explain, in one sentence each, which term of the bias--variance decomposition (Lecture 1, slide 41) is
responsible for each curve's behaviour.
\end{task}

\subsection{Learning curves}
\begin{task}{Task 7 --- how much data is enough?}
For training-set sizes $n\in\{20,40,80,160,320\}$, fit the linear model on the first $n$ training rows
and record training MSE and validation MSE (on a fixed 20\% validation slice of the training set).
Plot both against $n$. The gap between the curves is an estimate of the variance term; its closing rate
tells you whether more data would help this model.
\end{task}

\subsection{The one number}
Only now: evaluate the chosen model (degree from Task 6) on \texttt{X\_te}. Report the test MSE. Do not
go back and change anything afterwards.

\section{Exercises (to submit before Lab 2)}
\begin{exercise}
Implement mini-batch SGD for the least-squares loss with batch sizes $B\in\{1,16,442\}$ and a constant
$\eta$. Plot the loss after each epoch for 30 epochs, all three on one axis. Which one reaches the lowest
loss? Which one gets close fastest in wall-clock time? Explain both with Lecture 1, slides 35--36.
\end{exercise}
\begin{exercise}
Prove that for a random subset $\mathcal B$ of size $B$ drawn without replacement, the mini-batch
gradient $\frac1B\sum_{i\in\mathcal B}\nabla\ell_i(\vbeta)$ is an unbiased estimate of $\nabla L(\vbeta)$.
Then compute its variance in terms of $B$, $N$, and $\Var_i\bigl(\nabla\ell_i(\vbeta)\bigr)$, and check
that it vanishes at $B=N$.
\end{exercise}
\begin{exercise}
On the diabetes data, add the $\binom{10}{2}=45$ pairwise interaction features. Training $R^2$ goes up.
Does 5-fold CV $R^2$? Report both, and state what the result says about $N=442$ relative to $p=55$.
\end{exercise}

\section*{Submission}
A single notebook, all cells executed top to bottom without error (\emph{Kernel $\to$ Restart \& Run All}),
uploaded to the course drive before Wednesday, September 16, 18:00.\\ Name it \texttt{lab01\_lastname\_firstname.ipynb}.

\end{document}
